% page 281 of the facsimile (image p-280 of the scan)
% block 162.6 x 267.9 mm ; left margin 39.4 mm
\pgc{17.780mm}{0.000mm}{
\l{}
\l{}
\l{}
\l{}
\l{\cel{52}273}
\l{}
\l{kaulo. (\sou{bot.}) Planto supala ek la familio \textquotedbl{}kruciferi\textquotedbl{}. - L.}
\l{\cel{3}\sou{brassica oleracea}. - DL.}
\l{}
\l{kaurio. Konko uzata kom +koino en Bengal ed en la regioni cen-}
\l{\cel{3}trala di Afrika.\cel{3}L.\cel{1}\sou{cypraea}. - FSL}
\l{}
\l{kaustika. Qua desorganizas, korodas la tisui di la animali e}
\l{\cel{3}di la vejetanti. - DEFIS.}
\l{}
\l{\sou{kaustiko}. (\sou{geom.}) Kurvo qua envelopas lumo-radii refraktita}
\l{\cel{3}da spegulo konkava, ta radii venante de un lumo-punto.}
\l{\cel{3}- DEFIS.}
\l{}
\l{\sou{kautero}. Korpo varmega o kemiajo, quan onu uzas medicinale}
\l{\cel{3}por desorganizar tisuo. - DEFIS.}
\l{}
\l{\sou{kauzo}. I. (\sou{filoz}.)To per quo kozo esas o divenas to quon lu}
\l{\cel{3}esas. - II.To per quo ulo eventas ; to pro quo onu agas o}
\l{\cel{3}facas ulo. - DEFIS.}
\l{}
\l{\sou{kava}. Qua havas en su vakuo plu o min profunda.}
\l{}
\l{\sou{kavalo}. (\sou{zool}.) Animalo domestika, mamifera, ek la familio}
\l{\cel{3}\textquotedbl{}solipedi\textquotedbl{}, qua uzesas kom kavalkajo, kom tirbestio o kom}
\l{\cel{3}charjo-bestio. - L.\cel{1}\sou{Caballus}. DEFIRS.}
\l{}
\l{\sou{kavalkar}. (\sou{netrans.}) Sidar quale viro sur kavalo, sur biciklo.}
\l{\cel{3}- EFIRS.}
\l{}
\l{\sou{kavalrio}. La parto di armeo, qua konsistas ye viri sur kavali.}
\l{\cel{3}- DEFIRS.}
\l{}
\l{\sou{kavatino}. (\sou{muziko)}. Kant-ario (\textquotedbl{}moderato\textquotedbl{}) di la operi italiana.}
\l{\cel{3}- DEFIRS.}
\l{}
\l{\sou{kaverno}. Kavajo naturala,sub sulo o sub roko, ed apta kom}
\l{\cel{3}refujeyo. - EFIRS.}
\l{}
\l{\sou{kaviaro}. Disho acesora Rusa, ek sturg-ovi,forte presita e}
\l{\cel{3}marinigita. - DEFIRS.}
\l{}
\l{\sou{kayo}. I. Ter-amasajo, subtenata da muro ek quadro, e protektata}
\l{\cel{3}da parapeto, facita alonge rivero por restriktar lu dum}
\l{\cel{3}acenso, en la parto di lua fluo qua trairas urbo. - II. Rivo}
\l{\cel{3}di portuo, ube onu kargas o deskargas la vari. - III.}
\l{\cel{3}(\sou{metaf}.) kayo por la relvoyi. - DEF.}
\l{}
\l{\sou{kayero.}\cel{2}I. Asemblajo ek plura papero-folii, generale inter-}
\l{\cel{3}sutita unlatere. - II. Asemblajo ek la folii ye qui konsistas}
\l{\cel{3}fragmento di verko skribala quan onu editas partope. - F.}
\l{}
\l{kazo. I. Ek la eventaji quin povas produktar stando, konjunturo,}
\l{\cel{3}la eventajo qua esas produktita, fakte o supozate o konjek-}
\l{\cel{3}tate. II. (gram.) En ula lingui, dezinenco di la substantivo}
\l{\cel{3}(e di la adjektivo e pronomo qui relatas lu), por indikar lua}
\l{\cel{3}precipua roli en la frazo. - DEFIRS.}
}
